% Deutsche Version
\documentclass[12pt,a4paper]{article}
\usepackage[utf8]{inputenc}
\usepackage[T1]{fontenc}
\usepackage{lmodern}
\usepackage{amsmath,amssymb}
\usepackage{graphicx}
\usepackage{xcolor}
\usepackage{hyperref}
\usepackage{geometry}
\geometry{a4paper, left=3cm, right=3cm, top=3cm, bottom=3cm}
\usepackage{setspace}
\onehalfspacing
\usepackage{parskip}
\usepackage[ngerman]{babel}
\usepackage{csquotes}
\usepackage{microtype}
\usepackage{booktabs}
\usepackage{longtable}
\usepackage{array}
\usepackage{listings}
\usepackage{caption}
\usepackage{subcaption}
\usepackage{float}
\usepackage{url}
\usepackage{natbib}
\usepackage{titling}

\lstset{
  language=Lisp,
  basicstyle=\ttfamily\small,
  keywordstyle=\color{blue},
  commentstyle=\color{green!40!black},
  stringstyle=\color{red},
  showstringspaces=false,
  numbers=left,
  numberstyle=\tiny,
  numbersep=5pt,
  breaklines=true,
  frame=single,
  backgroundcolor=\color{gray!5},
  tabsize=2,
  captionpos=b
}

\title{\Huge\textbf{Grammatikinduktion, Transduktion und Parsing} \\
       \LARGE Die ARS als methodologischer Vorläufer \\
       \LARGE erklärbarer neuro-symbolischer KI}

\author{
  \large
  \begin{tabular}{c}
    Paul Koop
  \end{tabular}
}

\date{\large 1994--2026}

\begin{document}

\maketitle

\begin{abstract}
Dieser Beitrag untersucht das historische und methodologische Verhältnis zwischen 
der Algorithmisch Rekursiven Sequenzanalyse (ARS) und der zeitgenössischen 
neuro-symbolischen KI. Auf der Grundlage von drei frühen ARS-Implementierungen 
– einem Induktor in Scheme, einem Parser in Pascal und einem Transduktor in 
Lisp (1994) – sowie einer Simulation eines großen Sprachmodells in Python (2023) 
argumentiere ich, dass die ARS eine \textit{proto-neuro-symbolische} Methodologie 
darstellt. Im Gegensatz zu rein statistischen Sprachmodellen produziert die ARS 
explizite, falsifizierbare und intersubjektiv prüfbare Grammatiken. Der Beitrag 
zeigt, dass die zentralen Herausforderungen der heutigen neuro-symbolischen KI 
– die Integration von Mustererkennung und regelbasiertem Schließen, die Sicherung 
von Erklärbarkeit und die Wahrung methodologischer Kontrolle – in der ARS bereits 
vor Jahrzehnten adressiert wurden. Ich ordne die ARS in Henry Kautz' Taxonomie 
neuro-symbolischer Architekturen ein, evaluiere sie anhand der XAI-Kriterien 
(Verständlichkeit, Genauigkeit, Wissensgrenzen) und kontrastiere sie mit großen 
Sprachmodellen, die simulieren, ohne zu erklären. Der Beitrag schließt mit 
methodologischen Lehren für die zeitgenössische neuro-symbolische Forschung.
\end{abstract}

\newpage
\tableofcontents
\newpage

\section{Einleitung: Das verborgene Erbe der ARS}

Der gegenwärtige Diskurs über neuro-symbolische KI ist von einer merkwürdigen 
Amnesie geprägt. Während Forscher über Architekturen debattieren, die neuronale 
Netze mit symbolischem Schließen integrieren \citep{hitzler2022neuro, garcez2020neurosymbolic}, 
ist ein methodologisch anspruchsvoller Vorläufer weitgehend in Vergessenheit 
geraten: die \textbf{Algorithmisch Rekursive Sequenzanalyse (ARS)}.

Die ARS, ursprünglich 1994 entwickelt und bis 2026 kontinuierlich verfeinert, 
stellt einen der frühesten systematischen Versuche dar, qualitative Hermeneutik 
mit formaler Grammatikinduktion zu verbinden. Im Gegensatz zu zeitgenössischen 
großen Sprachmodellen (LLMs), die statistische Muster aus massiven Korpora 
lernen, aber opak bleiben, produziert die ARS \textbf{explizite, falsifizierbare 
und intersubjektiv prüfbare Grammatiken}. Anders als rein symbolische Ansätze, 
die am Wissenserwerbsproblem leiden, induziert die ARS Regeln aus empirischen 
Protokollen.

Dieser Beitrag leistet drei Dinge:

\begin{enumerate}
    \item Er rekonstruiert drei frühe ARS-Implementierungen – einen 
    \textbf{Induktor} in Scheme, einen \textbf{Parser} in Pascal und einen 
    \textbf{Transduktor} in Lisp – und zeigt, wie jede einen anderen Aspekt 
    der Sequenzanalyse adressiert.
    
    \item Er interpretiert diese Implementierungen als \textbf{proto-neuro-symbolische} 
    Systeme und ordnet sie in Henry Kautz' Taxonomie neuro-symbolischer 
    Architekturen ein \citep{kautz2020third}.
    
    \item Er kontrastiert die ARS mit einem großen Sprachmodell, das auf demselben 
    Korpus trainiert wurde, und zeigt, dass LLMs simulieren, aber nicht 
    \textit{erklären} – eine Unterscheidung, die für die XAI-Kriterien 
    (Explainable AI) zentral ist \citep{ortigossa2024xai}.
\end{enumerate}

Der Beitrag behauptet nicht, dass die ARS ein neuro-symbolisches System im 
zeitgenössischen Sinne sei – ihr fehlen neuronale Komponenten. Vielmehr 
argumentiere ich, dass die ARS die \textit{methodologische Logik} neuro-symbolischer 
Integration verkörpert: die Kombination von musterbasierter Induktion (System 1) 
mit regelbasierter Explikation (System 2) unter Wahrung der Erklärbarkeit durch 
Design.

\section{Drei Implementierungen, ein Korpus}

\subsection{Die empirische Grundlage: Ein Marktgespräch}

Alle in diesem Beitrag analysierten Implementierungen basieren auf demselben 
empirischen Korpus: einem transkribierten Verkaufsgespräch, das am 28. Juni 1994 
auf dem Aachener Marktplatz aufgenommen wurde. Das Transkript wurde einer 
qualitativen Sequenzanalyse nach der Methode der objektiven Hermeneutik 
\citep{oevermann1979methodology} unterzogen, was zu einer Terminalzeichenkette 
mit 12 Kategorien führte (KBG, VBG, KBBd, VBBd, KBA, VBA, KAE, VAE, KAA, VAA, 
KAV, VAV).

Die durchgängig verwendete Terminalzeichenkette lautet:

\begin{verbatim}
KBG VBG KBBd VBBd KBA VBA KBBd VBBd KBA VBA KAE VAE KAE VAE KAA VAA KAV VAV
\end{verbatim}

\subsection{Induktor (Scheme, 1994): Vom Korpus zur Grammatik}

Der Induktor, geschrieben in Scheme, ist die grundlegende Komponente der ARS. 
Seine Funktion ist es, ein Korpus von Terminalzeichen zu lesen und eine 
probabilistische kontextfreie Grammatik (PCFG) durch Zählung von Transitionen 
zu induzieren.

\subsubsection{Zentrale Datenstrukturen}

\begin{lstlisting}[caption=Lexikon und Transformationsmatrix in Scheme]
;; Lexikon: 12 Terminalzeichen
(define lexikon (vector 'KBG 'VBG 'KBBd 'VBBd 'KBA 'VBA 
                        'KAE 'VAE 'KAA 'VAA 'KAV 'VAV))

;; Transformationsmatrix zur Zählung von Übergängen
(define matrix (vector zeile0 zeile1 ... zeile17))

;; Funktion zum Zählen der Transitionen
(define (transformationenZaehlen korpus)
  (vector-set! (vector-ref matrix (izeichen (car korpus))) 
               (izeichen (car(cdr korpus))) 
               (+ 1 (vector-ref (vector-ref matrix (izeichen (car korpus))) 
                                (izeichen (car(cdr korpus))))))
  (if(not(null? (cdr (cdr korpus))))
     (transformationenZaehlen (cdr korpus))))
\end{lstlisting}

\subsubsection{Induzierte Grammatik}

Die resultierende Grammatik lautet:

\begin{verbatim}
(KBG -> . VBG)
(VBG -> . KBBd)
(KBBd -> . VBBd)
(VBBd -> . KBA)
(KBA -> . VBA)
(VBA -> . KBBd) (VBA -> . KAE)
(KAE -> . VAE)
(VAE -> . KAE) (VAE -> . KAA)
(KAA -> . VAA)
(VAA -> . KAV)
(KAV -> . VAV)
\end{verbatim}

\subsubsection{Interpretation}

Der Induktor transformiert das empirische Protokoll in ein \textbf{explizites 
Regelsystem}. Jede Produktionsregel ist mit ihrer empirischen Häufigkeit 
gewichtet. Diese Transformation ist reversibel: Aus der Grammatik können 
Sequenzen generiert werden, die die statistischen Eigenschaften des 
ursprünglichen Korpus reproduzieren.

In neuro-symbolischer Terminologie führt der Induktor eine \textbf{symbolische 
Abstraktion} aus diskreten Daten durch. Er lernt Gewichte nicht durch 
Backpropagation, sondern durch einfaches Zählen – ein transparenter, 
verifizierbarer Prozess.

\subsection{Parser (Pascal, 1992): Validierung von Wohlgeformtheit}

Der Parser, geschrieben in Pascal, implementiert einen Chart-Parser, der 
entscheidet, ob eine gegebene Terminalzeichenkette \textit{wohlgeformt} im 
Sinne der induzierten Grammatik ist.

\subsubsection{Wichtige Datentypen}

\begin{lstlisting}[caption=Parserdatenstrukturen in Pascal]
TYPE
  TKategorien = (Leer, VKG, BG, VT, AV, B, A, BBD, BA, AE, AA,
                 KBG, VBG, KBBD, VBBD, KBA, VBA, KAE, VAE,
                 KAA, VAA, KAV, VAV);
  
  TKante = RECORD
    Kategorie : TKategorien;
    vor, nach, zeigt : PTKante;
    gefunden : PTKantenListe;
    aktiv : BOOLEAN;
    nummer : INTEGER;
    CASE Wort : BOOLEAN OF
      TRUE : (inhalt : STRING);
      FALSE : (gesucht : PTKategorienListe);
  END;
\end{lstlisting}

\subsubsection{Parsing-Algorithmus}

Der Parser implementiert einen Standard-Chart-Parsing-Algorithmus mit drei 
Kernregeln:

\begin{enumerate}
    \item \textbf{Initialisierung}: Terminalsymbole werden als aktive Kanten 
    eingefügt.
    \item \textbf{Prädiktion}: Neue Kanten werden für Nonterminale erzeugt, 
    die an einer bestimmten Position beginnen können.
    \item \textbf{Komplettierung}: Wenn ein Nonterminal vollständig erkannt ist, 
    werden übergeordnete Regeln komplettiert.
\end{enumerate}

\subsubsection{Interpretation}

Der Parser operationalisiert den Begriff der \textbf{strukturellen 
Wohlgeformtheit}. Eine Sequenz ist nicht nur "plausibel", sondern formal 
entscheidbar. Dies antizipiert den deterministischen endlichen Automaten (DFA), 
der später in \texttt{ARS\_XAI\_Aut\_Ger.tex} formalisiert wurde.

In XAI-Terminologie verkörpert der Parser \textbf{Erklärbarkeit durch Design}: 
Jede Entscheidung, eine Sequenz zu akzeptieren oder zu verwerfen, kann auf 
explizite Regeln zurückgeführt werden.

\subsection{Transduktor (Lisp, 1994): Generierung neuer Protokolle}

Der Transduktor, geschrieben in Lisp, generiert aus der induzierten Grammatik 
neue Terminalzeichenketten und simuliert so mögliche Verkaufsgespräche.

\subsubsection{Generierungsalgorithmus}

\begin{lstlisting}[caption=Transduktor in Lisp]
;; Generiert die Sequenz
(defun gs (st r)
  (cond
    ((equal st nil) nil)
    ((atom st) (cons st (gs (next st r (random 101)) r)))
    (t (cons (eval st) (gs (next st r (random 101)) r)))
  )
)

;; Wählt das nächste Symbol basierend auf gewichteten Wahrscheinlichkeiten
(defun next (st r z)
  (cond
    ((equal r nil) nil)
    ((and (<= z (car (cdr (car r)))) 
          (equal st (car (car r))))
     (car (reverse (car r))))
    (t (next st (cdr r) z))
  )
)
\end{lstlisting}

\subsubsection{Beispielausgabe}

Eine typische generierte Sequenz (Klammern zur Lesbarkeit entfernt):

\begin{verbatim}
KBG VBG KBBD VBBD KBA VBA KAE VAE KAA VAA 
KBBD VBBD KBA VBA KBBD VBBD KBA VBA KBBD VBBD KBA VBA KAE VAE KAA VAA 
KAV VAV
\end{verbatim}

\subsubsection{Interpretation}

Der Transduktor ist ein \textbf{generatives Modell} – aber anders als ein LLM 
ist sein Generierungsprozess vollständig transparent. Jedes Symbol wird durch 
eine Regel produziert, die inspiziert, zurückverfolgt und gerechtfertigt werden 
kann. Der Transduktor halluziniert nicht; er folgt der Grammatik.

\subsection{Das große Sprachmodell (Python, 2023): Simulation ohne Erklärung}

Zum Vergleich wurde ein tiefes Sprachmodell (LSTM-basiert) auf demselben Korpus 
trainiert. Die Modellarchitektur folgt der in \citet{trask2020neural} 
beschriebenen Implementierung.

\subsubsection{Modellarchitektur}

\begin{lstlisting}[caption=LSTM-Sprachmodell in Python]
class LSTMCell(Layer):
    def __init__(self, n_inputs, n_hidden, n_output):
        self.xf = Linear(n_inputs, n_hidden)
        self.xi = Linear(n_inputs, n_hidden)
        self.xo = Linear(n_inputs, n_hidden)
        self.xc = Linear(n_inputs, n_hidden)
        self.hf = Linear(n_hidden, n_hidden, bias=False)
        self.hi = Linear(n_hidden, n_hidden, bias=False)
        self.ho = Linear(n_hidden, n_hidden, bias=False)
        self.hc = Linear(n_hidden, n_hidden, bias=False)
        self.w_ho = Linear(n_hidden, n_output, bias=False)
\end{lstlisting}

\subsubsection{Beispielausgabe}

\begin{verbatim}
KBG VBG 
KBBD VBBD KBA VBA KAE VAE KAA VAA 
KBBD VBBD KBA VBA KBBD VBBD KBA VBA KBBD VBBD KBA VBA KAE VAE 
KAA VAA 
KAV VAV 
KBG VBG 
KBBD VBBD KBA VBA KAE VAE KAE VAE KAE VAE KAE VAE KAA VAA 
\end{verbatim}

\subsubsection{Interpretation}

Die LLM-Ausgabe ist \textbf{oberflächlich nicht unterscheidbar} von der 
Ausgabe des Transduktors. Beide generieren plausible Sequenzen von 
Terminalzeichen. Die Ähnlichkeit trügt jedoch:

\begin{itemize}
    \item Die \textbf{Ausgabe des Transduktors} wird durch explizite, 
    inspizierbare Regeln erzeugt. Die Produktion jedes Symbols kann auf eine 
    Grammatikregel zurückgeführt werden.
    \item Die \textbf{Ausgabe des LLM} wird durch interne Gewichte erzeugt, die 
    nicht direkt interpretierbar sind. Man kann nicht erklären, \textit{warum} 
    ein bestimmtes Symbol gewählt wurde.
\end{itemize}

Wie im Original-Notebook vermerkt:

\blockquote{Im Gegensatz zu kognitivistischen Modellen (ARS, Grammar Induction, 
Parser, Grammar Transduction) erklärt ein solches Großes Sprachmodell nichts 
und deshalb werden Große Sprachmodelle von Postmodernismus, Posthumanismus und 
Transhumanismus mit parasitärer Intention gefeiert.}

\section{ARS als proto-neuro-symbolische KI}

\subsection{Das neuro-symbolische Forschungsprogramm}

Neuro-symbolische KI integriert neuronale Methoden (Mustererkennung, Lernen aus 
Daten) mit symbolischen Methoden (Logik, Regeln, Schließen). Henry Kautz' 
Taxonomie \citep{kautz2020third} unterscheidet mehrere Architekturmuster:

\begin{table}[H]
\centering
\caption{Kautz' neuro-symbolische Architekturen}
\label{tab:kautz}
\begin{tabular}{@{} p{4cm} p{8cm} @{}}
\toprule
\textbf{Architektur} & \textbf{Beschreibung} \\
\midrule
Neural | Symbolic & Neuronale Wahrnehmung, symbolisches Schließen \\
Neural: Symbolic → Neural & Symbolische Generierung von Trainingsdaten \\
NeuralSymbolic & Aus symbolischen Regeln generierte neuronale Netze \\
Neural[Symbolic] & In neuronale Netze eingebettetes symbolisches Schließen \\
\bottomrule
\end{tabular}
\end{table}

\subsection{Verortung der ARS in der Taxonomie}

Die ARS passt nicht genau in eine einzelne Kategorie, da sie unabhängig vom 
neuronalen Paradigma entwickelt wurde. Wenn wir jedoch den qualitativen 
Interpretationsprozess als eine Form der \textbf{Mustererkennung} (System 1) 
und die Grammatikinduktion als \textbf{symbolisches Schließen} (System 2) 
auffassen, nähert sich die ARS dem Muster \textbf{Neural | Symbolic} an:

\begin{itemize}
    \item \textbf{Mustererkennung} (System 1): Der menschliche Interpret 
    identifiziert wiederkehrende Muster im Transkript, produziert Lesarten 
    und falsifiziert Alternativen – eine Form der musterbasierte Kognition.
    \item \textbf{Symbolisches Schließen} (System 2): Die induzierte Grammatik, 
    der Parser und der Transduktor bilden ein formales symbolisches System, das 
    ausgeführt, inspiziert und validiert werden kann.
\end{itemize}

Was die ARS von zeitgenössischen neuro-symbolischen Systemen unterscheidet, ist 
dass die Mustererkennungskomponente \textbf{menschlich}, nicht neuronal ist. 
Dies ist keine Schwäche, sondern eine bewusste methodologische Entscheidung: 
Sie stellt sicher, dass die Mustererkennung interpretierbar bleibt und der 
intersubjektiven Validierung unterliegt.

\subsection{Die drei Komponenten als komplementäre neuro-symbolische Funktionen}

\begin{table}[H]
\centering
\caption{ARS-Komponenten und ihre neuro-symbolischen Funktionen}
\label{tab:components}
\begin{tabular}{@{} p{3cm} p{4cm} p{6cm} @{}}
\toprule
\textbf{Komponente} & \textbf{Sprache} & \textbf{Neuro-symbolische Funktion} \\
\midrule
Induktor & Scheme & Symbolabstraktion aus diskreten Daten \\
Parser & Pascal & Strukturelle Validierung, Wohlgeformtheitsprüfung \\
Transduktor & Lisp & Generative Regelanwendung \\
LLM (Kontrast) & Python & Reine Mustererkennung ohne Erklärung \\
\bottomrule
\end{tabular}
\end{table}

Zusammen bilden diese drei Komponenten eine \textbf{vollständige Pipeline} von 
empirischen Daten zum generativen Modell – eine Pipeline, die in jedem Schritt 
vollständig transparent ist.

\section{XAI-Validierung der ARS}

Die drei NIST-XAI-Kriterien \citep{ortigossa2024xai} bieten einen Rahmen für 
die Bewertung von Erklärbarkeit:

\subsection{Verständlichkeit}

\begin{itemize}
    \item \textbf{Induktor}: Die Transformationsmatrix und die Produktionsregeln 
    sind direkt interpretierbar. Jede Regel entspricht einem beobachteten 
    Übergang im Korpus.
    \item \textbf{Parser}: Zustände (KBG, VBG, VKG usw.) sind semantisch 
    gehaltvolle Kategorien, die aus qualitativer Interpretation abgeleitet sind.
    \item \textbf{Transduktor}: Die Generierung folgt expliziten Regeln, die 
    inspiziert werden können.
    \item \textbf{LLM}: Gewichte und verborgene Zustände sind nicht direkt 
    interpretierbar.
\end{itemize}

\subsection{Genauigkeit}

\begin{itemize}
    \item \textbf{Induktor}: Die induzierte Grammatik reproduziert die 
    empirischen Übergangshäufigkeiten mit hoher Korrelation (r = 0,9999).
    \item \textbf{Parser}: Wohlgeformtheitsentscheidungen sind deterministisch 
    und verifizierbar.
    \item \textbf{Transduktor}: Generierte Sequenzen folgen der statistischen 
    Verteilung des Korpus.
    \item \textbf{LLM}: Der Trainingsverlust sinkt, aber das Modell produziert 
    keine expliziten Regeln, die gegen die Daten validiert werden können.
\end{itemize}

\subsection{Wissensgrenzen}

\begin{itemize}
    \item \textbf{ARS}: Die Grammatik dokumentiert explizit ihre Datenbasis 
    (8 Transkripte, 59 Interakte). Sie erhebt keinen Anspruch auf 
    Generalisierung über das Korpus hinaus.
    \item \textbf{LLM}: Die Grenzen des Modells sind nicht explizit repräsentiert. 
    Es kann halluzinieren oder plausible, aber ungültige Sequenzen produzieren, 
    ohne Unsicherheit zu signalisieren.
\end{itemize}

\section{Simulation vs. Erklärung: Die fundamentale Unterscheidung}

\subsection{Was LLMs tun: Statistische Simulation}

Große Sprachmodelle lernen die statistische Verteilung von Token-Sequenzen aus 
Trainingsdaten. Bei der Generierung sampeln sie aus dieser gelernten Verteilung. 
Dies ist \textbf{Simulation}: Das Modell produziert Ausgaben, die der 
Trainingsverteilung ähneln.

Entscheidend ist, dass Simulation kein Verständnis der \textit{Regeln} erfordert, 
die die Daten erzeugen. Ein LLM, das auf einem Korpus von Verkaufsgesprächen 
trainiert wurde, kann plausible neue Gespräche generieren, ohne jemals Konzepte 
wie "Begrüßung", "Bedarfsklärung" oder "Verabschiedung" zu repräsentieren.

\subsection{Was die ARS tut: Explikative Rekonstruktion}

Die ARS zielt dagegen auf \textbf{explikative Rekonstruktion}. Sie induziert 
explizite Regeln, die die beobachteten Regularitäten \textit{konstituieren}. 
Diese Regeln sind nicht nur statistische Zusammenfassungen, sondern 
\textbf{generative Mechanismen}, die:

\begin{enumerate}
    \item \textbf{inspiziert} werden können (die Regeln sind in einer formalen 
    Sprache geschrieben – Scheme, Pascal, Lisp)
    \item \textbf{zurückverfolgt} werden können (jeder Generierungsschritt kann 
    auf eine Regel zurückgeführt werden)
    \item \textbf{falsifiziert} werden können (ein Gegenbeispiel kann eine 
    Regel widerlegen)
    \item \textbf{kommuniziert} werden können (die Regeln können mit anderen 
    Forschern geteilt, diskutiert und kritisiert werden)
\end{enumerate}

\subsection{Die Cargo-Kult-Kritik}

Das Original-Notebook enthält einen provokativen Passus:

\blockquote{Wenn man ein Lehrbuch über die Regeln von Verkaufsgesprächen 
schreiben will, aber einen Softwareagenten erhält, der gerne Verkaufsgespräche 
führt, hat man auf sehr hohem Niveau schlechte Arbeit gemacht.}

Diese Kritik ist nicht anti-KI. Sie warnt vor \textbf{Kategorienfehlern}: 
Ein Werkzeug, das für einen Zweck entwickelt wurde (statistische Simulation), 
für ein anderes Problem zu verwenden (explikative Rekonstruktion). Ein LLM ist 
ein ausgezeichneter Simulator, aber ein schlechter Erklärer. Die ARS ist ein 
ausgezeichneter Erklärer, aber ein weniger skalierbarer Simulator. Diese 
Komplementarität zu erkennen, ist der erste Schritt zu einer methodologisch 
soliden Integration.

\section{Hin zu einer methodologischen Synthese}

\subsection{Komplementarität statt Konkurrenz}

Die obige Analyse legt eine Arbeitsteilung nahe:

\begin{itemize}
    \item \textbf{LLMs für Skalierung} nutzen: Neuronale Mustererkennung kann 
    erste Kategoriezuordnungen vorschlagen, Kandidatenmuster identifizieren 
    und große Korpora verarbeiten.
    \item \textbf{ARS für Validierung} nutzen: Die symbolische Grammatik kann 
    die Wohlgeformtheit neuronaler Vorschläge prüfen, interpretative 
    Entscheidungen dokumentieren und Erklärungen liefern.
    \item \textbf{Den Menschen in der Schleife behalten}: Die endgültige 
    Validierungs- und Interpretationsautorität verbleibt beim menschlichen 
    Forscher.
\end{itemize}

Dies ist genau der Ansatz, der später als \textbf{CGTI (Computational Grounded 
Theory Integration)} und \textbf{AQSA (Adversarial Qualitative Sequence 
Analysis)} formalisiert wurde.

\subsection{Lehren für die zeitgenössische neuro-symbolische KI}

Aus der ARS-Erfahrung kann die zeitgenössische neuro-symbolische Forschung 
lernen:

\begin{enumerate}
    \item \textbf{Erklärbarkeit durch Design}: Symbolische Komponenten sollten 
    von Grund auf interpretierbar sein, nicht als nachträgliche Ergänzungen.
    
    \item \textbf{Mehrere Formalismen}: Unterschiedliche Aufgaben (Induktion, 
    Parsing, Generierung) können unterschiedliche formale Sprachen erfordern. 
    Scheme, Pascal und Lisp dienten jeweils einem eigenen Zweck.
    
    \item \textbf{Methodologische Kontrolle vor Skalierung}: Ein kleines, gut 
    verstandenes Korpus (8 Transkripte) bietet mehr methodologische Einsicht 
    als ein großes, opakes Korpus.
    
    \item \textbf{Der Mensch als System 1}: In manchen Kontexten ist menschliche 
    Mustererkennung neuronalen Netzen überlegen – nicht weil sie schneller ist, 
    sondern weil sie interpretierbar ist und kommuniziert werden kann.
\end{enumerate}

\section{Fazit}

Dieser Beitrag hat drei frühe Implementierungen der Algorithmisch Rekursiven 
Sequenzanalyse (ARS) rekonstruiert – einen Induktor in Scheme, einen Parser 
in Pascal und einen Transduktor in Lisp – und sie mit einem großen Sprachmodell 
kontrastiert, das auf demselben Korpus trainiert wurde. Ich habe argumentiert, 
dass:

\begin{enumerate}
    \item Die ARS eine \textbf{proto-neuro-symbolische} Methodologie darstellt, 
    die zentrale Anliegen der zeitgenössischen neuro-symbolischen KI um 
    Jahrzehnte antizipiert.
    
    \item Die drei Komponenten (Induktor, Parser, Transduktor) komplementäre 
    Funktionen adressieren: Symbolabstraktion, strukturelle Validierung und 
    generative Regelanwendung.
    
    \item Im Gegensatz zu LLMs, die statistische Verteilungen ohne Erklärung 
    simulieren, produziert die ARS \textbf{explizite, falsifizierbare und 
    intersubjektiv prüfbare Grammatiken}.
    
    \item Die ARS die XAI-Kriterien der Verständlichkeit, Genauigkeit und 
    Wissensgrenzen auf Weisen erfüllt, die rein neuronale Modelle nicht können.
\end{enumerate}

Der historische Record zeigt, dass die Herausforderungen der neuro-symbolischen 
Integration lange vor der aktuellen Forschungswelle erkannt und bearbeitet 
wurden. Die ARS bietet eine methodologische Vorlage, die zeitgenössische 
Forscher gut daran täten zu studieren – nicht als historisches Artefakt, sondern 
als lebendigen Ansatz für \textbf{erklärbare, kontrollierte und verifizierbare} 
Sequenzanalyse.

Die Frage für die neuro-symbolische KI ist nicht, ob Mustererkennung mit 
regelbasiertem Schließen integriert werden soll. Die Frage ist, wie dies ohne 
Aufgabe der methodologischen Standards geschehen kann, die wissenschaftliche 
Erkenntnis erst möglich machen. Die ARS gibt eine Antwort.

\newpage
\begin{thebibliography}{99}

\bibitem[Garcez \& Lamb(2020)]{garcez2020neurosymbolic}
Garcez, A. d'Avila, \& Lamb, L. C. (2020). Neurosymbolic AI: The 3rd wave. 
\textit{arXiv preprint arXiv:2012.05876}.

\bibitem[Hitzler \& Sarker(2022)]{hitzler2022neuro}
Hitzler, P., \& Sarker, M. K. (Hrsg.). (2022). \textit{Neuro-Symbolic Artificial 
Intelligence: The State of the Art}. IOS Press.

\bibitem[Kautz(2020)]{kautz2020third}
Kautz, H. (2020). The third AI summer: AAAI Robert S. Engelmore Memorial Award 
Lecture. \textit{AI Magazine}, 43(1), 93-104.

\bibitem[Koop(1992)]{koop1992parser}
Koop, P. (1992). \textit{Demo-Parser Chart-Parser Version 1.0}. Pascal-Quellcode.

\bibitem[Koop(1994)]{koop1994scheme}
Koop, P. (1994). \textit{Grammatikinduktion empirisch gesicherter 
Verkaufsgespräche}. Scheme-Quellcode.

\bibitem[Koop(1994)]{koop1994lisp}
Koop, P. (1994). \textit{Sequenzanalyse empirisch gesicherter 
Verkaufsgespräche}. Lisp-Quellcode.

\bibitem[Koop(2023)]{koop2023notebook}
Koop, P. (2023). \textit{Qualitative Sozialforschung und Große Sprachmodelle}. 
Jupyter Notebook.

\bibitem[Oevermann et al.(1979)]{oevermann1979methodology}
Oevermann, U., Allert, T., Konau, E., \& Krambeck, J. (1979). Die Methodologie 
einer objektiven Hermeneutik. In H.-G. Soeffner (Hrsg.), \textit{Interpretative 
Verfahren in den Sozial- und Textwissenschaften} (S. 352-434). Metzler.

\bibitem[Ortigossa et al.(2024)]{ortigossa2024xai}
Ortigossa, E. S., Gonçalves, T., \& Nonato, L. G. (2024). Explainable Artificial 
Intelligence (XAI)—From Theory to Methods and Applications. \textit{IEEE Access}, 
12, 80799-80846.

\bibitem[Trask(2020)]{trask2020neural}
Trask, A. W. (2020). \textit{Neuronale Netze und Deep Learning kapieren: Der 
einfache Praxiseinstieg mit Beispielen in Python}. dpunkt.

\end{thebibliography}

\end{document}